%% kapa: exam=1 round=36 year=2025 period=1 form=A subject=civil_law subject_ko=민법 seq=1 number=1
%% source: Q-Net 감정평가사 기출문제 (cst00309) artlSeq=5231525 file=raw/1008/5231525/2025년 제36회 감정평가사 1차시험 1교시 문제지 원본.pdf
%% conversion: kordoc 4.24.51 -> tools/gen1.py (auto); review=proofread-sonnet; flags=none
\begin{question}{1}
신의성실의 원칙에 관한 설명으로 옳지 않은 것은? (다툼이 있으면 판례에 따름)
\begin{choices}
  \item 채무자의 소멸시효를 이유로 한 항변권의 행사도 신의성실의 원칙의 지배를 받는다.
  \item 사정변경으로 인한 계약해제에 있어 ‘사정’에는 일방 당사자의 주관적인 사정은 포함되지 않는다.
  \item 강행규정을 위반한 법률행위를 한 자가 그 법률행위의 무효를 주장하는 것은 원칙적으로 신의칙에 위배된다.
  \item 법원은 신의성실의 원칙의 위반 여부를 직권으로 판단할 수 있다.
  \item 변호사의 민사소송위임사무에 관한 약정보수액이 부당하게 과다하여 신의성실의 원칙에 반하는 경우, 변호사는 적당하다고 인정되는 범위 내의 보수액만을 청구할 수 있다.
\end{choices}
\end{question}
\answer{3}
